\clearpage
\fancyhead[L]{\small FLORA.}
\SourceLeaf{481}
\begin{center}
\rule{0.55\linewidth}{0.4pt}\\[2pt]
Astronomische Nachrichten. Vol. XXVI. No. 622. p. 352. 10 January 1848.\\[2pt]
\rule{0.55\linewidth}{0.4pt}
\end{center}

I have received the following observations of \emph{Flora}, made with the Göttingen meridian circle, from Privy Councillor \textsc{Gauss}.

\begin{center}
\setlength{\tabcolsep}{4pt}
\begin{tabular}{l|r r r|r|r}
Dec. 10 & $11^{\rm u}$ & $11^{\rm m}$ & $4^{\rm s}0$ & $66^\circ\ 47'\ 26''2$ & $+14^\circ\ 22'\ 50''1$ \\
\hfill 11 & & $6$ & $7.9$ & $66\ 32\ 5.4$ & $14\ 25\ 42.6$ \\
\hfill 12 & & $1$ & $11.0$ & $66\ 16\ 59.9$ & $14\ 28\ 40.0$ \\
\end{tabular}
\end{center}

Privy Councillor \textsc{Gauss} notes that the R.A. on the 10th is somewhat less reliable than on the 11th, and that Flora appears much brighter than Iris. The observations could be made with full illumination, and the planet is a fine sight above $\rho$ Tauri, which, a few minutes farther north, followed it by $7^{\rm s}$.

